% ============================================================================
% CHEATSHEET TEMPLATE
% ============================================================================
% Based on the ANA-G cheatsheet style
%
% Features:
% - 4-column landscape layout on A4
% - Colored boxes for main concepts and sub-concepts
% - Compact formatting optimized for information density
% - Mathematical typesetting support
%
% Usage:
% 1. Replace metadata (title, author, date)
% 2. Adjust number of columns if needed (see \Ncolumn)
% 3. Add your content using mainbox and subbox environments
% 4. Compile with pdflatex (run twice for proper layout)
% ============================================================================

\documentclass[a4paper, fontsize=7pt]{scrartcl}

% ============================================================================
% ENCODING & LANGUAGE
% ============================================================================
\usepackage[utf8]{inputenc}
% \usepackage[german]{babel}    % Uncomment for German hyphenation
% \usepackage[english]{babel}   % Uncomment for English hyphenation

% ============================================================================
% PAGE LAYOUT
% ============================================================================
% 4-column landscape layout with minimal margins
\usepackage[landscape, left=0.75cm, top=1cm, right=0.75cm, bottom=1.5cm, footskip=5pt]{geometry}
\usepackage{flowfram}
\ffvadjustfalse
\setlength{\columnsep}{1cm}

% Number of columns (adjust as needed: 3-4 recommended for A4)
\Ncolumn{4}

% Remove paragraph indentation
\setlength{\parindent}{0cm}

% ============================================================================
% COLORS & BOXES
% ============================================================================
\usepackage[most]{tcolorbox}

% Base style for all boxes
\tcbset{
  base/.style={
    boxrule=0.3mm,
    leftrule=1mm,
    left=1.75mm,
    arc=0mm,
    fonttitle=\bfseries,
    colbacktitle=black!10!white,
    coltitle=black,
    toptitle=0.75mm,
    bottomtitle=0.25mm,
    title={#1}
  }
}

% Main concept box (green accent)
% Use for: definitions, theorems, important formulas
\definecolor{maincolor}{rgb}{0.004, 0.5, 0.35}  % Nice green
\newtcolorbox{mainbox}[1]{
  colframe=maincolor,
  base={#1}
}

% Sub concept box (gray accent)
% Use for: examples, remarks, minor concepts
\newtcolorbox{subbox}[1]{
  colframe=black!20!white,
  base={#1}
}

% Alternative colors (uncomment to use)
% \definecolor{maincolor}{rgb}{0.2, 0.4, 0.8}   % Blue
% \definecolor{maincolor}{rgb}{0.8, 0.2, 0.2}   % Red
% \definecolor{maincolor}{rgb}{0.5, 0.2, 0.6}   % Purple

% ============================================================================
% MATHEMATICS
% ============================================================================
\usepackage{amsmath, amsthm, mathtools, amssymb}
\usepackage{amsfonts, amstext}
\allowdisplaybreaks  % Allow page breaks in equations

% Common math shortcuts (customize as needed)
\def\limn{\lim_{n\to \infty}}
\def\limxo{\lim_{x\to 0}}
\def\limxi{\lim_{x\to\infty}}
\def\limxn{\lim_{x\to-\infty}}
\def\sumk{\sum_{k=1}^\infty}
\def\sumn{\sum_{n=0}^\infty}

% Number sets
\def\N{\mathbb{N}}  % Natural numbers
\def\Z{\mathbb{Z}}  % Integers
\def\Q{\mathbb{Q}}  % Rationals
\def\R{\mathbb{R}}  % Reals
\def\C{\mathbb{C}}  % Complex

% Calculus
\def\dx{\,dx}
\def\dt{\,dt}
\def\dy{\,dy}
\newcommand{\ddx}[1]{\frac{d#1}{dx}}
\newcommand{\ddn}[2]{\frac{d#1}{d#2}}
\newcommand{\pdn}[2]{\frac{\partial #1}{\partial #2}}

% Big O notation
\def\oof{\mathcal{O}}

% ============================================================================
% TABLES
% ============================================================================
\usepackage{tabularx, multirow, booktabs}
\renewcommand*{\arraystretch}{1.5}  % More spacing in tables

% ============================================================================
% LISTS
% ============================================================================
\usepackage{enumitem}
\setitemize{itemsep=0.5pt, topsep=2pt}
\setenumerate{itemsep=0.75pt, topsep=2pt}

% ============================================================================
% GRAPHICS & TIKZ
% ============================================================================
\usepackage{graphicx}
\usepackage{tikz}
\usetikzlibrary{
  positioning,
  calc,
  fit,
  matrix,
  decorations,
  angles,
  quotes,
  babel,
  3d,
  arrows.meta,
  shadings,
  shapes.geometric
}

\tikzset{
  >=stealth,
  % Add custom TikZ styles here
}

% ============================================================================
% HYPERLINKS
% ============================================================================
\usepackage{hyperref}
\hypersetup{
  colorlinks=true,
  linkcolor=maincolor,
  urlcolor=maincolor
}

% ============================================================================
% ADDITIONAL PACKAGES
% ============================================================================
\usepackage{centernot}   % For \centernot command
\usepackage{marvosym}    % Additional symbols
\usepackage{wasysym}     % Additional symbols

% ============================================================================
% DOCUMENT METADATA
% ============================================================================
\title{Your Cheatsheet Title}
\author{Your Name}
\date{\today}

% Table of contents depth
\setcounter{tocdepth}{4}
\setcounter{secnumdepth}{4}

% ============================================================================
% DOCUMENT START
% ============================================================================
\begin{document}

\maketitle

% ============================================================================
% SECTION 1: EXAMPLE SECTION
% ============================================================================
\section{First Topic}

\begin{mainbox}{Main Concept Title}
This is a \textbf{main box} for important concepts, definitions, or theorems.

\textbf{Key Formula:}
$$\boxed{E = mc^2}$$

Use the boxed command to highlight the most important formulas.
\end{mainbox}

\begin{subbox}{Example or Remark}
This is a \textbf{sub box} for examples, remarks, or secondary information.

\textbf{Example:} Calculate $E$ for $m = 1$ kg:
$$E = 1 \cdot (3 \times 10^8)^2 = 9 \times 10^{16} \text{ J}$$
\end{subbox}

% ============================================================================
% SUBSECTION EXAMPLE
% ============================================================================
\subsection{Subsection Title}

\begin{mainbox}{Another Main Concept}
\textbf{Definition:} A clear definition goes here.

\textbf{Properties:}
\begin{itemize}
    \item First property
    \item Second property
    \item Third property
\end{itemize}
\end{mainbox}

\begin{subbox}{Worked Example}
\textbf{Problem:} Find the derivative of $f(x) = x^2$.

\textbf{Solution:}
$$f'(x) = \lim_{h \to 0} \frac{(x+h)^2 - x^2}{h} = \lim_{h \to 0} \frac{2xh + h^2}{h} = 2x$$
\end{subbox}

% ============================================================================
% SECTION 2: TABLES EXAMPLE
% ============================================================================
\section{Tables and Lists}

\begin{mainbox}{Comparison Table}
\renewcommand*{\arraystretch}{1.4}
\begin{center}
\begin{tabular}{l|c|c}
\textbf{Method} & \textbf{Pros} & \textbf{Cons} \\
\hline
Method A & Fast & Less accurate \\
Method B & Accurate & Slower \\
Method C & Balanced & Complex \\
\end{tabular}
\end{center}
\end{mainbox}

\begin{subbox}{Enumerated Steps}
\textbf{Algorithm:}
\begin{enumerate}
    \item Initialize variables
    \item Iterate until convergence
    \item Return result
\end{enumerate}
\end{subbox}

% ============================================================================
% SECTION 3: FORMULAS REFERENCE
% ============================================================================
\section{Formula Reference}

\begin{mainbox}{Important Formulas}
\textbf{Derivatives:}
\begin{align*}
\ddx{x^n} &= nx^{n-1} \\
\ddx{e^x} &= e^x \\
\ddx{\ln x} &= \frac{1}{x}
\end{align*}
\end{mainbox}

\begin{mainbox}{Integrals}
\textbf{Basic Integrals:}
\begin{align*}
\int x^n \dx &= \frac{x^{n+1}}{n+1} + C \\
\int e^x \dx &= e^x + C \\
\int \frac{1}{x} \dx &= \ln|x| + C
\end{align*}
\end{mainbox}

% ============================================================================
% SECTION 4: REFERENCE TABLES
% ============================================================================
\section{Quick Reference}

\begin{mainbox}{Symbol Reference}
\renewcommand*{\arraystretch}{1.3}
\begin{center}
\begin{tabularx}{\linewidth}{lX}
$\N$ & Natural numbers $\{1, 2, 3, \ldots\}$ \\
$\Z$ & Integers $\{\ldots, -1, 0, 1, \ldots\}$ \\
$\R$ & Real numbers \\
$\C$ & Complex numbers \\
\end{tabularx}
\end{center}
\end{mainbox}

% ============================================================================
% TIPS FOR AVOIDING OVERFULL BOXES
% ============================================================================
%
% 1. Keep formulas concise - split long equations over multiple lines
% 2. Use shorter variable names in inline math
% 3. Abbreviate table headers
% 4. Use align* environment for multi-line equations
% 5. Test compile frequently to catch overflow early
%
% Example of splitting a long formula:
% Instead of: $$a + b + c + d + e = f + g + h + i + j$$
% Use:
% \begin{align*}
% a + b + c + d + e &= f + g + h \\
%                   &\quad + i + j
% \end{align*}
%
% ============================================================================

\end{document}